% !TEX program = xelatex

%% main.tex
%% Copyright 2025 Guodong Li
%
% This work may be distributed and/or modified under the
% conditions of the LaTeX Project Public License, either version 1.3
% of this license or (at your option) any later version.
% The latest version of this license is in
%   https://www.latex-project.org/lppl.txt
% and version 1.3c or later is part of all distributions of LaTeX
% version 2008 or later.
%
% This work has the LPPL maintenance status `maintained'.
% 
% The Current Maintainer of this work is Guodong Li.
%
% This work consists of the files sdu-thesis.cls main.tex

\documentclass{settings/sdu-thesis}

% 论文信息
% !TEX root = ../main.tex

% info.tex
% 中图分类号
\fenlei{TP309}

% 单位代号
\DWdaihao{10422}
% 密级
\miji{公\quad 开}
% 学号
\StuNum{123456789}
% 页眉
\ThesisHeader{山东大学博士学位论文}
\EThesisHeader{Dissertation for Doctoral Degree}
% 论文类型（封面） 中
\CThesisType{博~士~学~位~论~文}
% 论文类型（封面） 英
\EThesisType{Dissertation for Doctoral Degree}
% 中文标题
\Ctitle{山东大学博士学位论文模板}
% 英文标题
\Etitle{A Template for Doctoral Dissertation of Shandong University}
% 作者中文名
\Cauthor{哈利波特}
% 学院
\Dpart{格兰芬多学院}
% 专业中文名
\Cmajor{魔法专业}
% 年级
\Grade{}
% 指导老师中文名
\Csuperver{笨笨虎钓鱼}
\Ccosuperver{}
% 中文日期
\Cdate{2025年1月4日}
%\Cdate{\today}
\title{\the\Ctitle}
\author{\the\Cauthor}
\date{\the\Cdate}

% 学位类型
\DegreeType{}        %学术学位 不填此项
%\DegreeType{（专业学位）} %专业学位

\begin{document}

\maketitlepagestatement

% 中文目录 和 英文目录
\content

% !TEX root = ../main.tex

\begin{cnabstract}

    本文介绍了一个用于撰写山东大学博士学位论文的 \LaTeX{} 模板。该模板旨在帮助博士研究生按照学校的
    格式要求，快速、高效地编写和排版学位论文。模板包括了封面、摘要、目录、正文、参考文献等部分的格式
    设置，并提供了详细的使用说明和示例代码。通过使用该模板，用户可以专注于论文内容的撰写，而无需担心
    排版问题，从而提高论文写作的效率和质量。

    本模板是在 sduthesis.cls(\url{https://github.com/Liam0205/sduthesis/}) 这一模板的基
    础上进行修改的，感谢原作者的辛勤劳动。该模板在参考的模板的基础上，修改了部份格式，消除了部份关于
    字体的警告，使之能在我的环境下正常编译。然而，由于本人技术水平有限，模板难免存在一些问题和不足之
    处，还请各位使用者批评指正，共同完善。

    \cnkeywords{山东大学； 博士学位论文； \LaTeX{} 模板}

\end{cnabstract}

\begin{enabstract}

    This document introduces a \LaTeX{} template for writing doctoral dissertations
    at Shandong University. The template is designed to help doctoral students
    write and format their dissertations quickly and efficiently according to the
    school's requirements. The template includes format settings for the cover,
    abstract, table of contents, main text, references, etc., and provides detailed
    instructions and sample code. By using this template, users can focus on
    writing the content of the dissertation without worrying about formatting
    issues, thereby improving the efficiency and quality of dissertation writing.

    This template is based on the sduthesis.cls template
    (\url{https://github.com/Liam0205/sduthesis/}), and I would like to thank the
    original author for his hard work. This template modifies some formats based on
    the referenced template, eliminates some warnings about fonts, and allows it to
    compile normally in my environment. However, due to my limited technical level,
    the template inevitably has some problems and shortcomings. I hope that users
    can criticize and correct them to improve them together.

    \enkeywords{Shandong University; Doctoral Thesis; \LaTeX{} Template}

\end{enabstract}

% 插图目录 和 表格目录
\ftcontent

\chapter{模板使用示例}
\echapter{Template Usage Example}
\section{模板下载}
\esection{Template Download}
你可以从以下网址下载本模板：\url{https://github.com/gudngli/SDU-Thesis}。

\section{基本结构}
\esection{Basic Structure}
本模板的基本结构如下：
\begin{itemize}
    \item settings/：模板设置文件夹，包含论文的格式设置。
    \item main.tex：主文件，包含论文的整体结构。
    \item chapters/：章节文件夹，包含各个章节的内容。
    \item ref.bib：参考文献文件，包含所有引用的文献信息。使用cite命令引用参考文献，例如\cite{MSR-TIT24}。
\end{itemize}

\section{编译方法}
\esection{Compilation Method}
使用以下命令编译论文：
    \begin{verbatim}
        xelatex main.tex
        bibtex main.aux
        xelatex main.tex
        xelatex main.tex
    \end{verbatim}
或者
\begin{itemize}
    \item Windows：使用 TeX Live 或 MiKTeX 编译。
    \item MacOS：使用 MacTeX 编译。(或者终端运行``make''命令)
\end{itemize}
本文档已经在 TeXlive 2023 上测试通过。

\section{选项说明}
\esection{Options Description}

\begin{itemize}
    \item \texttt{[double]}: 双面打印。此时所有章节的首页都在奇数页，另外，页面边距也会有所调整，内侧边距比外侧边距大。以便装订。
    \item \texttt{[fzfont]}: 使用方正字体。默认使用的是宋体，如需在论文首页使用方正字体，请添加此选项。使用前确保自己的电脑上已经安装了``方正粗黑宋简体''字体。
    该字体的安装文件可以在settings文件夹下找到。
    \item \texttt{[print]}: 使用该选项时所有链接的颜色都会被设置为黑色。
\end{itemize}

\section{论文标题，作者信息设置}
\esection{Thesis Title, Author Information Setting}
在 \texttt{./chapters/info.tex} 文件中，你可以设置论文的标题，作者信息等。

\section{插入图片}
\esection{Insert Picture}
插入图片的示例代码如下：
\begin{verbatim}
    \begin{figure}[h]
        \centering
        \includegraphics[width=0.5\linewidth]{settings/logo.pdf}
        \caption[SDUlogo]{山东大学校徽}
        \label{fig:sdulogo}
    \end{figure}
\end{verbatim}
效果如下，请留意插图目录的变化：
\begin{figure}[h]
    \centering
    \includegraphics[width=0.5\linewidth]{settings/logo.pdf}
    \caption[SDUlogo]{山东大学校徽}
    \label{fig:sdulogo}
\end{figure}

\section{插入表格}
\esection{Insert Table}
插入表格的示例代码如下：
\begin{verbatim}
    \begin{table}[h]
        \centering
        \label{tab:sdulogo}
        \begin{tabular}{|c|c|}
            \hline
            山东大学 & 校徽 \\
            \hline
            山东大学 & 校徽 \\
            \hline
        \end{tabular}
        \caption{山东大学校徽}
    \end{table}
\end{verbatim}
效果如下，请留意表格目录的变化：
\begin{table}[h]
    \centering
    \label{tab:sdulogo}
    \begin{tabular}{|c|c|}
        \hline
        山东大学 & 校徽 \\
        \hline
        山东大学 & 校徽 \\
        \hline
    \end{tabular}
    \caption{山东大学校徽}
\end{table}

一个稍微复杂的表格示例：
\begin{table}[ht]
    \centering
    \renewcommand{\arraystretch}{1}
    \begin{tabular}{|l|l|c|}
        \hline
        \multicolumn{3}{|c|}{Optimal-access}                                                                       \\
        \hline
        \hline
        Codes                      & Sub-packetization $\ell$                                    & Restrictions    \\
        \hline
        \cite[Section VIII]{Ye16}  & $s^{n-1}$                                                   &                 \\
        \hline
        \cite{CB}                  & $s^{n}$                                                     &                 \\
        \hline
        \cite{Liu22}               & $s^{\left\lceil n/2 \right\rceil}$                          &                 \\
        \hline
        \cite{Vajha21}             & $s^{\left\lceil n/s \right\rceil}$      (smallest possible) & $s\in\{2,3,4\}$ \\
        \hline
        %Section~\ref{sect:code1}   & $s^{\left\lceil n/s \right\rceil}$      (smallest possible) &                 \\
        %\hline
        \multicolumn{3}{c}{}                                                                                       \\
        \hline
        \multicolumn{3}{|c|}{Not optimal-access}                                                                   \\
        \hline
        \hline
        Codes                      & Sub-packetization $\ell$                                    & Restrictions    \\
        \hline
        \cite[Section IV]{Ye16}    & $s^{n}$                                                     &                 \\
        \hline
        \cite{LWHY22}              & $s^{\left\lceil n/3 \right\rceil}$      (best-known)        & $d=k+1$         \\
        \hline
        \cite{Zhang23} \cite{Li23} & $s^{\left\lceil n/2 \right\rceil}$                          &                 \\
        \hline
        %Section~\ref{sect:code2}   & $s^{\left\lceil n/(s+1) \right\rceil}$    (best-known)      &                 \\
        %\hline
        %\multicolumn{3}{c}{}
    \end{tabular}
    \caption[MSR码的显式构造]{Explicit $(n,k)$ MSR code constructions with linear field size and $d<n-1$ helper nodes, where $s=d-k+1$.
        ``smallest possible'' means that $\ell$ achieves the theoretical lower bound; ``best-known'' means that $\ell$ is the smallest among all the existing constructions.}
    \label{tab:cmp}
\end{table}

% 参考文献
\references
\bibliographystyle{plain}
\bibliography{ref}

% 致谢
% !TEX root = ../main.tex

\acknowledgment

在此，我要特别感谢 \LaTeX{} 模板 sduthesis.cls 的原作者 Liam Huang（\url{https://github.com/Liam0205/sduthesis/}），
感谢他辛勤的劳动和无私的分享。正是因为有了他的工作，我才能在此基础上进行修改， 得到当前的 \LaTeX{} 模板。

感谢所有为开源社区做出贡献的开发者们，你们的努力使得学术研究和论文写作变得更加高效和便捷。

% 博士期间所发论文
% !TEX root = ../main.tex

\publications

\begin{enumerate}
      \item
            \textbf{Guodong Li}, Ningning Wang, Sihuang Hu, and Min Ye.
            \newblock {MSR Codes With Linear Field Size and Smallest Sub-Packetization for Any Number of Helper Nodes}.
            \newblock {\em IEEE Transactions on Information Theory}, pages 1--1, 2024.
\end{enumerate}

% 博士期间所获奖项
% !TEX root = ../main.tex

\awards

\begin{enumerate}
    \item 山东大学新生入学学业奖学金，2020.10
\end{enumerate}

% 学位论文评阅及答辩情况表
%\comments
%\begin{figure*}[htp]
%	\centering
%	\includegraphics[scale=0.82]{./contents/table.eps}
%\end{figure*}

\end{document}
